% ============================================================
% Sección 4: Agente vs Agente LLM
% ============================================================
\section{Agente clásico vs. agente LLM}

\begin{frame}{¿En qué se parecen?}
  \begin{itemize}
    \item Ambos siguen el ciclo \textbf{percibir $\to$ decidir $\to$ actuar}.
    \item Ambos operan dentro de un entorno (juego, sistema operativo,
      conversación, internet).
    \item Ambos pueden mejorar con retroalimentación.
  \end{itemize}
  \vfill
  \centering
  \Large ¿Entonces qué cambia?
\end{frame}

\begin{frame}{Diagrama comparativo}
  \centering
  \begin{tikzpicture}[node distance=0.9cm]
    % Agente clasico
    \node[cajaBase, fill=colGris!15, minimum width=4.6cm] (titulo1) {\textbf{Agente clásico (RL)}};
    \node[cajaAgente, below=of titulo1, minimum width=4.6cm] (pol) {Política $\pi(a|s)$\\aprendida por prueba/error};
    \node[cajaEntorno, below=of pol, minimum width=4.6cm] (acc) {Acciones discretas\\predefinidas por el entorno};

    % Agente LLM
    \node[cajaBase, fill=colGris!15, right=2.2cm of titulo1, minimum width=4.6cm] (titulo2) {\textbf{Agente LLM}};
    \node[cajaLLM, below=of titulo2, minimum width=4.6cm] (razon) {Razonamiento en lenguaje\\natural (chain-of-thought)};
    \node[cajaHerramienta, below=of razon, minimum width=4.6cm] (tools) {Acciones = llamadas a\\herramientas/funciones};

    \draw[flecha] (titulo1) -- (pol);
    \draw[flecha] (pol) -- (acc);
    \draw[flecha] (titulo2) -- (razon);
    \draw[flecha] (razon) -- (tools);
  \end{tikzpicture}
\end{frame}

\begin{frame}{Tabla comparativa}
  \small
  \begin{center}
  \begin{tabular}{@{}p{3.1cm}p{4.1cm}p{4.1cm}@{}}
    \toprule
    & \textbf{Agente clásico (RL)} & \textbf{Agente LLM} \\
    \midrule
    Cómo decide & Política numérica aprendida & Razonamiento en lenguaje (prompt) \\
    Cómo aprende & Prueba y error con recompensa & Preentrenado en texto + fine-tuning \\
    Espacio de acciones & Fijo, definido por el entorno & Abierto: cualquier herramienta con schema \\
    Adaptación a tareas nuevas & Requiere reentrenar & Cambia con el \emph{prompt}, sin reentrenar \\
    Explicabilidad & Baja (pesos numéricos) & Alta (razona en texto legible) \\
    \bottomrule
  \end{tabular}
  \end{center}
\end{frame}

\begin{frame}{Entonces, ¿qué es un agente LLM?}
  \definicion{Agente LLM}{
    Un sistema en el que un modelo de lenguaje \textbf{decide qué hacer a
    continuación} —incluyendo qué herramientas invocar y con qué
    argumentos— de forma iterativa, usando su propio razonamiento en
    lenguaje natural como ``memoria de trabajo'', hasta cumplir un objetivo.
  }
  \vfill
  \begin{itemize}
    \item El LLM cumple el rol de \textbf{política}: dado un estado
      (contexto), produce una acción (texto, tool call).
    \item Las \textbf{herramientas} amplían el entorno con el que puede
      interactuar (buscar, ejecutar código, leer archivos...).
  \end{itemize}
\end{frame}
